\chapter{Accessibility and inclusive design}
\label{ch:a11y}

\begin{aside}
A TUI is often the most accessible kind of UI: screen readers
work with it, low-vision users can configure terminal fonts and
contrast, and there is no mouse-only flow. But ``most accessible
by default'' is not ``always accessible''. A few choices in your
app design can welcome or exclude users with disabilities. This
chapter walks the considerations and the patterns.
\end{aside}

\section{Why TUIs start accessible}

A typical TUI:

\begin{itemize}[leftmargin=*]
  \item Renders text, which screen readers can read.
  \item Uses high-contrast cells (the user picks the colours).
  \item Has keyboard navigation by default.
  \item Avoids animation that distracts.
  \item Works with the user's chosen font size.
\end{itemize}

These are wins by default. The patterns to maintain them
require thought.

\section{Screen readers}

Modern screen readers (Orca on Linux, JAWS / NVDA on Windows,
VoiceOver on macOS) can read terminal output, but they only see
what is on the screen \emph{after} ANSI escapes are interpreted.

A few principles:

\begin{itemize}[leftmargin=*]
  \item Information conveyed by colour alone is invisible to
    screen reader users. Pair it with a text marker.
  \item Status indicators that flash (using SGR 5 blink) are
    not announced. Use static markers.
  \item Boxes drawn with Unicode borders read as random
    characters. Provide an option to hide borders or use the
    \code{ASCII} border style.
\end{itemize}

\subsection*{Live regions}

When the UI updates rapidly (a log streaming, a chat coming
in), screen readers can drop messages. Patterns:

\begin{itemize}[leftmargin=*]
  \item Only announce new lines (no re-reads of the whole
    panel).
  \item Provide a ``pause output'' toggle so the user can
    catch up.
  \item Offer a transcript view that does not auto-scroll.
\end{itemize}

\section{Colour and contrast}

\code{NO\_COLOR=1} is a hint, not a request. Honour it
strictly:

\begin{lstlisting}
if (std::getenv(\"NO_COLOR\")) {
    cfg.theme = themes::Monochrome();
}
\end{lstlisting}

The monochrome theme uses only attributes (Bold, Underline,
Dim) for distinction. Information that was conveyed by colour
is still distinguishable.

\subsection*{Contrast ratios}

WCAG 2.1 defines contrast ratios:

\begin{itemize}[leftmargin=*]
  \item AA Normal text: 4.5:1.
  \item AA Large text: 3:1.
  \item AAA Normal text: 7:1.
\end{itemize}

For a TUI, ``large text'' rarely applies (cells are uniform).
Aim for 4.5:1 minimum on default themes; provide a
high-contrast option at 7:1 for users who need it.

\subsection*{Colour blindness}

About 8\% of men and 0.5\% of women have some form of colour
blindness. Common patterns:

\begin{itemize}[leftmargin=*]
  \item Do not rely on red vs green alone. Pair with shapes
    (\code{[+]} vs \code{[-]}).
  \item Use distinct hues, not adjacent ones (red + blue, not
    red + orange).
  \item Test palettes in a simulator (Color Oracle, Sim Daltonism).
\end{itemize}

\section{Keyboard navigation}

A TUI cannot rely on a mouse; mouse support is a convenience.
The keyboard must reach every interactive element.

\subsection*{Tab order}

Document the tab order explicitly. The visible order should
match the logical order:

\begin{itemize}[leftmargin=*]
  \item Top to bottom.
  \item Left to right (or right to left in RTL languages).
  \item Skip non-interactive elements.
\end{itemize}

\maya{}'s \code{FocusManager} (Chapter~\ref{ch:focus}) stores
order explicitly. Set it to match the visual order; do not
rely on registration order alone.

\subsection*{Shortcut redundancy}

For every shortcut, ask: can the user discover it? Patterns:

\begin{itemize}[leftmargin=*]
  \item A help overlay listing all bindings.
  \item Status bar hints (e.g. \code{[q] quit  [/] search}).
  \item A command palette (Ctrl-P) so commands are discoverable
    by name even if the user does not know the binding.
\end{itemize}

\section{Animation and motion}

Animation can:

\begin{itemize}[leftmargin=*]
  \item Distract or trigger discomfort for users with motion
    sensitivity.
  \item Mask important information.
  \item Slow down task completion.
\end{itemize}

Honour \code{prefers-reduced-motion}:

\begin{lstlisting}
if (std::getenv(\"MAYA_REDUCED_MOTION\")) {
    cfg.animations_enabled = false;
}
\end{lstlisting}

When disabled, spinners become static markers, smooth scrolling
becomes instant, transitions are skipped.

\subsection*{Avoid blinking}

SGR 5 (blink) was deprecated in many terminals and is generally
considered hostile. Do not use it.

\section{Font size and readability}

The user controls font size at the terminal level. Your app
should:

\begin{itemize}[leftmargin=*]
  \item Adapt to any terminal width (no hard-coded 80 columns).
  \item Wrap or truncate gracefully.
  \item Avoid tiny indicator characters (a single \code{x}
    when a clearer indicator like \code{[done]} is available).
\end{itemize}

\section{Time-sensitive UI}

If something requires the user to act quickly:

\begin{itemize}[leftmargin=*]
  \item Default timeouts should be generous.
  \item Provide a way to extend.
  \item Show the remaining time clearly.
\end{itemize}

Patterns like a confirmation prompt that closes after 5 seconds
are hostile to users who read slowly.

\section{Error messages}

Error messages should:

\begin{itemize}[leftmargin=*]
  \item Be visible to screen readers (not just colour).
  \item Persist long enough to read.
  \item Suggest a remedy (\emph{press r to retry}).
  \item Not blame the user.
\end{itemize}

A pattern: errors in a status bar with a fixed duration plus a
\code{[?]} indicator that opens a detail panel.

\section{Cognitive load}

Reduce things to remember:

\begin{itemize}[leftmargin=*]
  \item Consistent terminology across screens.
  \item Visible state (don't hide what mode you are in).
  \item Confirmations for destructive actions.
  \item Undo for everything reversible.
\end{itemize}

The Elm architecture's explicit \code{Model} is a gift here:
the state is a value the user can be shown.

\section{Internationalisation revisited}

Translation (Chapter~\ref{ch:i18n}) is an accessibility
concern: users whose first language is not the app's default
need translations.

The pattern: wrap all UX strings; ship at least one alternate
locale to prove the wrapping works.

\section{Test with assistive tech}

A few practical tests:

\begin{itemize}[leftmargin=*]
  \item Try the app with a screen reader.
  \item Try with no colour (\code{NO\_COLOR=1}).
  \item Try with a 16-colour terminal.
  \item Try with a 40-column-wide terminal.
  \item Try with motion disabled.
\end{itemize}

If any of these break the app, the app is not accessible.

\section{Pitfalls}

\begin{warning}
\textbf{Colour-only status indicators.} The most common
mistake: a red dot to mean ``error''. Add text or a shape so
non-colour users can see it.
\end{warning}

\begin{warning}
\textbf{Hard-coded narrow widths.} Apps that assume 80 columns
or 120 columns break at other sizes. \maya{}'s flex layout
helps; do not undo it with \code{width<80>}.
\end{warning}

\begin{warning}
\textbf{Rapid auto-updates.} A log scrolling past faster than
the user can read is hostile. Provide a pause toggle.
\end{warning}

\begin{warning}
\textbf{Modal dialogs that grab focus.} A modal that prevents
\code{Ctrl-C} from quitting traps the user. Always allow the
quit hotkey, even inside modals.
\end{warning}

\section*{Workshop}

\begin{enumerate}[leftmargin=*]
  \item Run your app with \code{NO\_COLOR=1}. Identify every
    place colour-only conveyed meaning. Add text or shapes.
  \item Test the app with Orca (Linux) or NVDA (Windows).
    Note which controls are not announced.
  \item Add a help overlay that lists every keyboard shortcut.
    Trigger it on \code{F1} or \code{?}.
\end{enumerate}

\begin{recap}
TUIs start accessible by default. Maintain that with: text
markers paired with colour, monochrome theme on NO\_COLOR,
explicit tab order, discoverable shortcuts, generous timeouts,
motion controls, persistent and respectful error messages.
Test with assistive tech. The Elm architecture's explicit
state makes accessibility easier, not harder.
\end{recap}
